\section{Mathematical framework}\label{sec:math}

\subsection{Semilinear problems and \texorpdfstring{$\varphi$}{phi} functions}

We consider finite-dimensional semidiscretizations of the form
\begin{equation}
  y'(t)=Ly(t)+N(t,y(t)),\qquad y(t_0)=y_0,
  \label{eq:semilinear}
\end{equation}
where $L$ is a linear operator and $N$ contains the remaining dynamics.  The
implementation does not require $L$ to be stored as a dense matrix: depending
on the backend, it may be diagonal data, a dense or sparse matrix, or an object
that supplies matrix-vector products.

The analytic coefficient functions are expressed through
\begin{equation}
 \varphi_0(z)=e^z,
 \qquad
 \varphi_k(z)=\int_0^1 e^{(1-\theta)z}
               \frac{\theta^{k-1}}{(k-1)!}\,d\theta,
 \quad k\ge1.
 \label{eq:phi}
\end{equation}
A method never needs a materialized matrix $\varphi_k(hL)$ as part of the
public execution contract.  It needs the \emph{action} of an analytic
coefficient, often a linear combination of several $\varphi_k(c hL)$, on one
or more vectors.  This distinction is central for large-scale problems and
for comparing Krylov and polynomial realizations.

\subsection{One explicit exponential GLM representation}

Let $X_n=(X_{n,1},\ldots,X_{n,r})^T$ denote the external state at time $t_n$.
The first component is the physical numerical solution.  For the historical
multistep methods represented in \pkg{}, the remaining components store scaled
nonlinear history.  The representation is the exponential general linear form
of Berland, Skaflestad and Wright~\cite{BerlandSkaflestadWright2007}, with the
order theory of Ostermann, Thalhammer and
Wright~\cite{OstermannThalhammerWright2006} and the B-series order conditions
of Berland, Owren and Skaflestad~\cite{BerlandOwrenSkaflestad2005} behind the
metadata; exponential Runge--Kutta schemes such as those of Cox and
Matthews~\cite{CoxMatthews2002} and Krogstad~\cite{Krogstad2005}, Lawson's
integrating-factor methods~\cite{Lawson1967} and the exponential Adams methods
of N{\o}rsett and of Hochbruck and Ostermann~\cite{HochbruckOstermann2011}
are its special cases.  A \texttt{MethodSpec} contains nodes $c$, analytic
operator-coefficient arrays $U,A,V,B$, and order/provenance metadata.  With
\begin{equation}
 G_j=hN(t_n+c_jh,Y_j),
\end{equation}
the common explicit engine evaluates
\begin{align}
 Y_i
   &=\sum_{j=1}^{r}U_{ij}(hL)X_{n,j}
     +\sum_{j<i}A_{ij}(hL)G_j,                        \label{eq:eglm-stage}\\
 X_{n+1,i}
   &=\sum_{j=1}^{r}V_{ij}(hL)X_{n,j}
     +\sum_{j=1}^{s}B_{ij}(hL)G_j.                   \label{eq:eglm-output}
\end{align}
The strict lower-triangular condition on $A$ encodes explicit stage ordering.
For $r>1$ the catalogue uses a nonlinear-history layout: after the physical
output row, the engine inserts $hN(t_n,y_n)$ and shifts older history.  This
makes multistep and general-linear schemes an external-state specialization of
the same step function rather than separate solver families.

The representation is deliberately descriptive.  Besides the coefficient
arrays, a method records its family, classical order, stiff order when known,
weak or conditional stiff order when distinct, stage order, quadrature order,
number of nonlinear evaluations, references, historical filename, notes and a
provenance status.  Consequently an executable method need not collapse
classical, stiff and weak-order statements into one integer.

\subsection{Order metadata are not interchangeable}

For nonstiff problems, the usual classical convergence order describes the
asymptotic temporal error under appropriate smoothness.  In semilinear stiff
problems, stronger operator-valued order conditions can determine whether the
same order is retained uniformly; failure of those conditions can produce
order reduction~\cite{HochbruckOstermann2005,HochbruckOstermann2010}.  Several
historical affine Lie-group entries additionally satisfy useful conditions at
$z=0$ without satisfying the corresponding strong identity for general $z$.
\pkg{} stores these as separate strong and weak or conditional metadata
rather than resolving a textual discrepancy by choosing the larger label.

The distinction matters experimentally.  A finite-grid manufactured problem
may show a higher apparent slope because a leading defect vanishes for that
problem.  Such a slope is useful validation evidence but does not by itself
change the catalogue's algebraic order metadata.  Section~\ref{sec:experiments}
therefore displays both ordinary classical tests and semidiscrete stiff-order
diagnostics.

\subsection{Analytic operator expressions}

Method coefficients are represented as analytic expression trees rather than
backend-specific arrays.  The basic object
\begin{equation}
  C(z)=\sum_{\ell} \alpha_\ell\varphi_{k_\ell}(c_\ell z)
  \label{eq:phi-combination}
\end{equation}
is a \texttt{PhiCombination}.  Sums of such expressions are canonicalized so
that a backend with native combination support can construct one numerical
action rather than evaluate every term independently.  Products represent
composition of analytic operator coefficients.  A shifted-resolvent expression
is kept separate from the $\varphi$ algebra and is used only by the retained
historical Crank--Nicolson comparator.

This symbolic layer is small but consequential.  It allows the catalogue to
state \emph{what} operator function a coefficient represents, while the backend
decides \emph{how} its action is approximated.  It is also the point at which
algebraic fusion becomes a semantics-level transformation rather than a
method-specific optimization.
